% ATTEMPT_STATUS: unresolved
% RESEARCH_GENERATED: 2026-10-01; GPT-6 Astra Ultra; individual mathematical attempt
% TOP_PROBLEM: {"id": 105, "title": "Linear Equation x + 3y = 2z + 2w", "category_id": 2, "category": {"id": 2, "name": "combinatorics", "display_name": "Combinatorics", "description": "Counting problems, graph theory, discrete structures.", "slug": "combinatorics", "order_index": 2, "created_at": "2026-07-31T15:26:25.671Z"}, "status": "open"}
% FIRST_POSTED: 2026-10-01
% Imported from: unsolvedmath_additional_500.tex; UnsolvedMath ID 105
% !TeX program = lualatex
% Compile twice: lualatex 105.tex
% Self-contained: no external bibliography, images, or \input files.
% Numeric section labels and visible numbers are original UnsolvedMath IDs.
\documentclass[11pt,a4paper]{article}
\usepackage[margin=24mm,headheight=14pt]{geometry}
\usepackage{fontspec}
\setmainfont{Latin Modern Roman}
\setsansfont{Latin Modern Sans}
\setmonofont{Latin Modern Mono}
\usepackage{amsmath,amssymb,mathtools}
\usepackage{unicode-math}
\setmathfont{Latin Modern Math}
\usepackage{microtype}
\usepackage{enumitem,xurl,needspace}
\usepackage[dvipsnames]{xcolor}
\usepackage{hyperref}
\hypersetup{unicode=true,colorlinks=true,linkcolor=MidnightBlue,urlcolor=MidnightBlue,
 pdftitle={UnsolvedMath Problem 105},pdfauthor={Source-annotated compilation},bookmarksnumbered=true}
\usepackage{bookmark,tocloft,titlesec,fancyhdr}
\setlength{\cftsecnumwidth}{4.2em}
\cftsetpnumwidth{2.5em}
\cftsetrmarg{4em}
\titleformat{\section}{\Large\bfseries\raggedright}{\thesection}{0.7em}{}
\pagestyle{fancy}\fancyhf{}
\fancyhead[L]{\small\sffamily UnsolvedMath research attempt}
\fancyhead[R]{\small Original catalogue IDs}
\fancyfoot[C]{\thepage}
\setlength{\parindent}{0pt}
\setlength{\parskip}{5pt plus 1pt}
\setlength{\emergencystretch}{3em}
\setcounter{tocdepth}{1}\setcounter{secnumdepth}{1}
\setlist[itemize]{leftmargin=3em,itemsep=4pt,topsep=4pt,parsep=0pt}
\allowdisplaybreaks
\newcommand{\N}{\mathbb{N}}
\newcommand{\Z}{\mathbb{Z}}
\newcommand{\Q}{\mathbb{Q}}
\newcommand{\R}{\mathbb{R}}
\newcommand{\C}{\mathbb{C}}
\newcommand{\F}{\mathbb{F}}
\newcommand{\A}{\mathbb{A}}
\newcommand{\PP}{\mathbb{P}}
\newcommand{\E}{\mathbb{E}}
\newcommand{\eps}{\varepsilon}
\newcommand{\dd}{\,\mathrm d}
\newcommand{\sref}[2]{\hyperlink{source:#1:#2}{[S#2]}}
\newif\ifshowcatalogue
\showcataloguetrue
\newcommand{\cataloguescope}[1]{\ifshowcatalogue
 \paragraph{Statement recorded in the supplied source.}
 #1\fi}
\begin{document}
% UnsolvedMath ID 105; source catalogue code GREEN-016

\setcounter{section}{104}

\section{Avoiding x + 3y = 2z + 2w}\label{105}

{\small\sffamily UnsolvedMath ID 105 · Combinatorics}

\subsection{Definitions and mathematical statement}

\paragraph{Shared notation.}Unless a section says otherwise, $\N=\{1,2,\ldots\}$,
$\Z$ is the ring of integers, $\Q,\R,\C$ are the rational, real, and complex
fields, $[N]=\{1,\ldots,\lfloor N\rfloor\}$, and $\log$ is the natural
logarithm. For real functions of a parameter tending to infinity,
$f\ll g$ and $f=O(g)$ mean $|f|\leq Cg$ eventually for some fixed $C$;
$f\gg g$ reverses this comparison; $f=o(g)$ means $f/g\to0$;
$f\sim g$ means $f/g\to1$; and $f\asymp g$ means both order bounds.
A subscript on $O$ or $\ll$ specifies allowed constant dependence.
The expression $n^{o(1)}$ in an upper bound means growth smaller than
$n^\epsilon$ for each fixed $\epsilon>0$. Local definitions govern any
exception, finite-field convention, or use of the same symbol for another
object. An asymptotic claim is not automatically an effective algorithm.



Define $F(N)$ as the maximum cardinality of $A\subseteq[N]$ for which there is no ordered quadruple $(x,y,z,w)\in A^4$ with all four entries pairwise distinct and $x+3y=2z+2w$. Repetitions are excluded only by that stated distinctness requirement. Determine the growth, or a sharp asymptotic description, of $F(N)$. \sref{105}{1} \sref{105}{2}

\paragraph{Statement recorded in the supplied source.}
What is the largest subset of $[N]$ with no solution to $x + 3y = 2z + 2w$ in distinct integers $x, y, z, w$? \sref{105}{1}

\subsection{Short English statement}

How many integers can be selected from an initial interval while avoiding this four-variable equation with four different selected values?

\subsection{Research attempt}

\paragraph{Provenance and scope.}
This individual attempt was generated by GPT-6 Astra Ultra on 1 October 2026. The method was to derive explicit partial results and test the first proposed extension adversarially. The original statement and sources below are retained. No claim of mathematical novelty or complete solution is made.

\paragraph{Formal target and context.}
Let $F(N)$ be the largest size of a subset of $[N]$ with no solution of $x+3y=2z+2w$ in four pairwise distinct elements. The equation is translation-invariant, since its coefficients sum to zero. Requiring all solutions to have $x=y=z=w$ is a stronger avoidance condition and is legitimate for a lower construction. The primary problem list attributes this equation to Ruzsa. The construction below is an elementary reconstruction of a quadratic obstruction, with no claim of novelty; its purpose is to produce a checked square-root-scale family and test the common sphere strategy.

\paragraph{A parabola with no nonconstant solutions.}
Let $p>3$ be prime such that $3$ is a quadratic nonresidue modulo $p$. In $\F_p^2$ take $P=\{(t,t^2):t\in\F_p\}$. Suppose four parameters $a,b,c,d$ satisfy both coordinate equations
\[
 a+3b=2c+2d,\qquad a^2+3b^2=2c^2+2d^2.
\]
Because two and four are invertible, write $a=b+4u$, $c=b+u+v$, and $d=b+u-v$. The first equation gives exactly this parametrisation. Substitution into the second gives
\[
 4b^2+8bu+16u^2=4b^2+8bu+4u^2+4v^2,
 \qquad v^2=3u^2.
\]
If $u\ne0$, division by $u^2$ makes $3$ a square, a contradiction. Hence $u=0$, then $v=0$, and all parameters coincide. Thus the parabola excludes even repeated-variable solutions other than the constant ones.

\paragraph{A carry-free embedding into an interval.}
For $0\leq t<p$, let $r_t$ be the least nonnegative residue of $t^2$ modulo $p$, set $B=4p$, and define $a_t=t+Br_t$. These are distinct because their residues modulo $B$ are distinct. If
\[
 a_x+3a_y=2a_z+2a_w,
\]
the lower-digit sums on both sides lie between zero and $4(p-1)<B$. Equality of the integers therefore implies equality of the lower digits and equality of the upper digits, with no carry ambiguity. Reducing those two equalities modulo $p$ gives the parabola equations just analysed. Thus the integer solution is constant. Since $0\leq a_t<4p^2$, translating by one places the set in $[4p^2]$ without changing the equation. We obtain $F(4p^2)\geq p$ for every such prime. The displayed construction and proof are unconditional at these parameters; no assertion about prime gaps is needed for this statement.

\paragraph{Why a sphere by itself does not provide near-linear examples.}
For three-term progressions, strict convexity can force three equal-norm points satisfying a midpoint relation to coincide. The present coefficients do not give that rigidity. Consider the four integer vectors in $\Z^4$
\[
 X=(-2,0,0,0),\quad Y=(2,0,0,0),\quad
 Z=(1,1,1,1),\quad W=(1,-1,-1,-1).
\]
They are pairwise distinct, all have squared norm four, and satisfy $X+3Y=2Z+2W$. Translating all four by $(2,2,2,2)$ makes every coordinate nonnegative while preserving the equation and their common sphere about the translated centre. Encoding the translated vectors in any base greater than sixteen gives four distinct integers on the same digit-sphere layer satisfying the equation without carries. Therefore the simple assertion that a fixed sphere layer automatically avoids the equation is false. Additional restrictions on the selected layer would be essential.

The comparison with the parabola also explains the role of the field hypothesis: over a field where three is a square, the equation $v^2=3u^2$ has nonzero solutions, so that particular construction loses its exclusion property. It is not enough that the curve is strictly convex over the real numbers; the arithmetic of the coefficient three is doing the work.

\paragraph{Verification and exact remaining gap.}
The companion script builds the digit sets for $p=5,7,17,19$, checks that three is nonsquare, and enumerates every ordered quadruple. The total solution counts are respectively $5,7,17,19$, exactly the constant solutions. The zero element in the computational sets is harmless because the final translation moves the construction into the positive interval. The proof treats repeated variables explicitly, so it certainly meets the weaker pairwise-distinct exclusion of the target. The result supplies a square-root-scale lower construction and an explicit obstruction to a naive sphere improvement. It does not show that square-root growth is optimal, nor approach a matching upper bound. The first remaining gap is a mechanism either to build much larger sets while controlling $v^2=3u^2$-type relations, or to force a genuinely four-distinct-variable solution above the conjectured extremal scale.

\subsection{Sources}

{\small

\begin{itemize}

\item[\hypertarget{source:105:1}{S1}] ULAM.AI, \emph{UnsolvedMath}, supplied \texttt{problems.json}, record 105 (GREEN-016), statement and background. The embedded literature-review date is 2026-08-17; that is the archive’s review date, not a fresh certification. Dataset landing page: \url{https://huggingface.co/datasets/ulamai/UnsolvedMath}.

\item[\hypertarget{source:105:2}{S2}] Tomasz Schoen and Olof Sisask, Roth's theorem for four variables and additive structures in sums of sparse sets, arXiv:1408.2568; Forum Math. Sigma 4 (2016), e5. \url{https://arxiv.org/abs/1408.2568}. Bibliographic information imported from the supplied record.

\item[\hypertarget{source:105:3}{S3}] Ben Green, 100 Open Problems, Problem 16, . \url{https://people.maths.ox.ac.uk/greenbj/papers/open-problems.pdf}. The author-hosted document was inspected on 27 September 2026; the archive’s local code is not a problem-number locator in this paper.

\end{itemize}

}

\label{end:105}
\end{document}
